\documentclass[10pt,a4paper]{ULBreport}
\usepackage[utf8]{inputenc}
\sceau{pic/official_logos/sceauULB.png}
\graphicspath{ {./pic/} }
\usepackage{multirow}
\usepackage{listings}
\usepackage{color} 
\usepackage{setspace} 
\usepackage{amsmath}
\usepackage[mathscr]{eucal}
\usepackage{hyperref}
\usepackage{pdfpages}
\usepackage{biblatex}
\usepackage{floatrow}
\usepackage{subcaption} 
\usepackage{siunitx}
\usepackage[many]{tcolorbox}
\usepackage{multirow}
\usepackage{listings}
\usepackage[dvipsnames]{xcolor}
\usepackage{fancyvrb}

\usepackage{xstring}
\usepackage{etoolbox}

% Colors



\begin{document} 


	\titleULB{
	title={V2V communication project},
    studies={M1-IRELE},
    course ={ELEC-H415 Communication Channels},
    author={\textit{Authors :} \\ Colot Emmeran },
    date={\textbf{Academic year :} \\ 2024 - 2025},
    teacher={\textit{Professor : } \\ Horlin François},
    logo={pic/official_logos/logos.jpg},
    manyAuthor
	}

%\listoftables % ToC for tables

%\listoffigures % ToC for figures

\chapter{Introduction}

\chapter{Theoretical answers}

\section{Step 1}

\begin{figure}[H]
    \centering
    \includegraphics[width=1\textwidth]{circuit.png}
    \caption{Equivalent electric circuit of an antenna in RX (left) and TX (right)}
    \label{fig:equivalent_electrical_circuit}
\end{figure}

Fig \ref{fig:equivalent_electrical_circuit} shows the equivalent electrical circuit at RX and TX where $\underline{V}_{oc}$ is the induced voltage, $\underline{V}_{RX}$ the voltage at the output of the RX antenna, $\underline{V}_{TX}$ at the input of the TX antenna and $\underline{I}_{a}$ the current entering the TX antenna.\\

As both transmitting and receiving antenna are vertical $\lambda/2$ dipoles, their equivalent heights can be analytically computed:

\begin{align*}
    \vec{h}_e (\theta, \phi) = \frac{\lambda}{\pi} \frac{\cos(\frac{1}{2}\cos \theta)}{\sin ^2 \theta}\vec{1_z}\\
    \vec{h}_{e\perp} (\theta, \phi) = -\frac{\lambda}{\pi} \frac{\cos(\frac{\pi}{2}\cos \theta)}{\sin \theta}\vec{1_\theta}
\end{align*}

The transverse part of the equivalent height allows to give an expression for the emitted electric field.

\begin{align*}
    \underline{\vec{E}} = j\frac{Z_0\underline{I}_a}{2\pi}\frac{\cos(\frac{\pi}{2}\cos \theta)}{\sin \theta}\frac{e^{-j\beta r}}{r}\vec{1_\theta}
\end{align*}

To make the transmission parameters appear in the electric field expression, $\beta$ and $\underline{I}_a$ must be replaced using:

\begin{align*}
    \beta = \frac{2\pi f_c}{c}\\
    \underline{V}_{TX} = (Z_a + Z_S) \underline{I}_a
\end{align*}

Wich yields

\begin{align*}
    \underline{\vec{E}} = j\frac{Z_0\underline{V}_{TX}}{2\pi (Z_a + Z_S)}\frac{\cos(\frac{\pi}{2}\cos \theta)}{\sin \theta}\frac{e^{-j \frac{2\pi f_c r}{c}}}{r}\vec{1_\theta}\\
    \left . \underline{\vec{E}}\right\vert_{\theta = \pi/2} = -j\frac{Z_0\underline{V}_{TX}}{2\pi (Z_a + Z_S)}\frac{e^{-j \frac{2\pi f_c r}{c}}}{r}\vec{1_z}
\end{align*}

Where $\theta$ has been chosen equal to $\pi /2$ in order to study only the electric filed in the horizontal place. At the receiving side, the incoming electric filed $\underline{\vec{E}}_i$ is simply equal to the transmitted electric field after it traveled over a distance $r = c\tau$ where $\tau$ is the propagation delay. This is of course in the case of line of sight (LOS) communication with no interfering object (IO) between TX and RX.

\begin{align*}
    \underline{\vec{E}}_i = -j\frac{Z_0 c \underline{V}_{TX}}{2\pi (Z_a + Z_S)}\frac{e^{-j 2\pi f_c \tau}}{\tau}\vec{1_z}
\end{align*}

The voltage at the output of the antenna $\underline{V}_{RX}$ can then be deduced using:

\begin{align*}
    \underline{V}_{RX} = \frac{Z_L}{Z_a + Z_L} \underline{V}_{oc}\\
    \underline{V}_{oc} = -\left . \vec{h}_{e\perp}(\theta, \phi)\right\vert_{\theta = \pi/2} \cdot \underline{\vec{E}}_i
\end{align*}

Resulting in:

\begin{align*}
    \underline{V}_{oc} &= \frac{\lambda}{\pi} \cdot \underline{E}_i\\
    &= -j\frac{Z_0 c \lambda \underline{V}_{TX}}{2\pi^2 (Z_a + Z_S)}\frac{e^{-j 2\pi f_c \tau}}{\tau}\\
    \underline{V}_{RX} &= -j\frac{Z_0 Z_L c \lambda \underline{V}_{TX}}{2\pi^2 (Z_a + Z_S) (Z_a+Z_S)}\frac{e^{-j 2\pi f_c \tau}}{\tau}
\end{align*}

\huge Is the actual formula correct?


%\printglossary

%\printglossary[type=\acronymtype]

%Bibliography
%\nocite{*}
%\printbibliography[type=article,title=Articles]

\end{document}	
